
\begin{center}
  \textbf{\large{LỜI CẢM ƠN}}
\end{center}
\addcontentsline{toc}{chapter}{LỜI CẢM ƠN}

Trước tiên, tôi xin bày tỏ lòng biết ơn sâu sắc nhất đến Giáo viên hướng dẫn của tôi,
TS. Vương Thị Hồng, người đã tận tâm định hướng, kiên nhẫn đồng hành và góp ý chỉnh
sửa chi tiết cho từng chương của khóa luận này. Sự nghiêm khắc khoa học của cô trong
việc buộc tôi phải ``phân biệt giữa một kết quả thực tế và một câu chuyện kể về kết
quả đó'' đã trở thành nguyên tắc cốt lõi chi phối toàn bộ phương pháp luận của công
trình nghiên cứu này.

Tôi xin chân thành cảm ơn các thầy cô giảng viên và cán bộ Khoa Công nghệ thông tin,
Trường Đại học Công nghệ, Đại học Quốc gia Hà Nội, đã trang bị cho tôi nền tảng kiến
thức vững chắc và kỷ luật kỹ thuật cần thiết qua bốn năm đại học. Đó là hành trang
không thể thiếu để tôi hoàn thành dự án kỹ thuật phức tạp này.

Tôi cũng gửi lời cảm ơn chân thành đến các bạn bè và đồng nghiệp đã tham gia vào
thí nghiệm đánh giá con người (Axis C) của công trình; sự chú tâm, kiên nhẫn và trung
thực của các bạn trong việc xếp hạng các bản tóm tắt cố tình được thiết kế để khó phân
biệt đã mang lại giá trị khoa học thực sự cho trục đánh giá thứ ba của khóa luận.

Cuối cùng, tôi xin gửi lời tri ân sâu sắc đến gia đình, những người đã bao dung,
thấu hiểu và kiên nhẫn chịu đựng những khoảng lặng dài trong những tháng ngày dự
án chưa đem lại kết quả như kỳ vọng, và là những người đầu tiên chúc mừng khi
kết quả cuối cùng đã đến.

\vspace{0.6em}
\begin{flushright}
  \textit{Hà Nội, tháng 5 năm 2026}\\
  \textbf{Lê Văn Thắng}
\end{flushright}
